\documentclass[UTF8,11pt]{ctexart}
\usepackage[a4paper,margin=24mm]{geometry}
\usepackage{amsmath,amssymb,amsthm,mathtools}
\usepackage[all]{xy}
\usepackage{xurl,hyperref,enumitem}
\hypersetup{hidelinks,breaklinks=true}
\setlength{\emergencystretch}{3em}
\newtheorem*{theorem}{Theorem}
\newtheorem*{lemma}{Lemma}
\newtheorem*{proposition}{Proposition}
\newtheorem*{corollary}{Corollary}
\theoremstyle{definition}
\newtheorem*{definition}{Definition}
\newtheorem*{remark}{Remark}
\DeclareMathOperator{\Hom}{Hom}
\DeclareMathOperator{\Ext}{Ext}
\DeclareMathOperator{\Tor}{Tor}
\DeclareMathOperator{\Spec}{Spec}
\DeclareMathOperator{\supp}{supp}
\DeclareMathOperator{\im}{im}
\DeclareMathOperator{\id}{id}
\DeclareMathOperator{\Tr}{Tr}
\DeclareMathOperator{\rank}{rank}
\newcommand{\R}{\mathbb R}
\newcommand{\C}{\mathbb C}
\newcommand{\Z}{\mathbb Z}
\newcommand{\N}{\mathbb N}
\newcommand{\Q}{\mathbb Q}
\begin{document}
\section*{014\quad 群行列式的不可约因子与表示维数}
\subsection*{来源与收录范围}
Pavel Etingof, Oleg Golberg, Sebastian Hensel, Tiankai Liu, Alex Schwendner, Dmitry Vaintrob, Elena Yudovina, \emph{Introduction to Representation Theory}；MIT OpenCourseWare 18.712, Fall 2010。使用课程所附 Chapter 4, \emph{Representations of finite groups: further results} 的 PDF，不使用另行出版的 AMS 图书。
\par\url{https://ocw.mit.edu/courses/18-712-introduction-to-representation-theory-fall-2010/84358595a02a73bced2c4e363a5d66f0_MIT18_712F10_ch4.pdf}
\par\url{https://ocw.mit.edu/courses/18-712-introduction-to-representation-theory-fall-2010/pages/lecture-notes/}
\par 获取日期：2026-09-28。\par 位置：\S4.2，Definition 4.6、Theorem 4.7、Lemma 4.8 及完整证明，分章 PDF 第2--3页。
\subsection*{作者命题与人类证明}

Enumerate the elements of a finite group $G$ as follows: $g_1,g_2,\ldots,g_n$. Introduce $n$ variables indexed with the elements of $G$:
\[x_{g_1},x_{g_2},\ldots,x_{g_n}.\]
\begin{definition}[4.6]
Consider the matrix $X_G$ with entries $a_{ij}=x_{g_i g_j}$. The determinant of $X_G$ is some polynomial of degree $n$ of $x_{g_1},x_{g_2},\ldots,x_{g_n}$ that is called the Frobenius determinant.
\end{definition}
\begin{theorem}[4.7]
\[
\det X_G=\prod_{j=1}^r P_j(x)^{\deg P_j}
\]
for some pairwise non-proportional irreducible polynomials $P_j(x)$, where $r$ is the number of conjugacy classes of $G$.
\end{theorem}
We will need the following simple lemma.
\begin{lemma}[4.8]
Let $Y$ be an $n\times n$ matrix with entries $y_{ij}$. Then $\det Y$ is an irreducible polynomial of $\{y_{ij}\}$.
\end{lemma}
\begin{proof}
Let $Y=t\cdot Id+\sum_{i=1}^n x_i E_{i,i+1}$, where $i+1$ is computed modulo $n$, and $E_{i,j}$ are the elementary matrices. Then $\det(Y)=t^n-(-1)^n x_1\cdots x_n$, which is obviously irreducible. Hence $\det(Y)$ is irreducible (since factors of a homogeneous polynomial are homogeneous).
\end{proof}
Now we are ready to proceed to the proof of Theorem 4.7.
\begin{proof}
Let $V=\C[G]$ be the regular representation of $G$. Consider the operator-valued polynomial
\[L(x)=\sum_{g\in G}x_g\rho(g),\]
where $\rho(g)\in\operatorname{End}V$ is induced by $g$. The action of $L(x)$ on an element $h\in G$ is
\[
L(x)h=\sum_{g\in G}x_g\rho(g)h=\sum_{g\in G}x_g gh
=\sum_{z\in G}x_{zh^{-1}}z.
\]
So the matrix of the linear operator $L(x)$ in the basis $g_1,g_2,\ldots,g_n$ is $X_G$ with permuted columns and hence has the same determinant up to sign.

Further, by Maschke's theorem, we have
\[\det_V L(x)=\prod_{i=1}^r\bigl(\det_{V_i}L(x)\bigr)^{\dim V_i},\]
where $V_i$ are the irreducible representations of $G$. We set $P_i=\det_{V_i}L(x)$. Let $\{e_{im}\}$ be bases of $V_i$ and $E_{i,jk}\in\operatorname{End}V_i$ be the matrix units in these bases. Then $\{E_{i,jk}\}$ is a basis of $\C[G]$ and
\[L(x)|_{V_i}=\sum_{j,k}y_{i,jk}E_{i,jk},\]
where $y_{i,jk}$ are new coordinates on $\C[G]$ related to $x_g$ by a linear transformation. Then
\[P_i(x)=\det|_{V_i}L(x)=\det(y_{i,jk}).\]
Hence, $P_i$ are irreducible (by Lemma 4.8) and not proportional to each other (as they depend on different collections of variables $y_{i,jk}$). The theorem is proved.
\end{proof}

\subsection*{整理说明（非作者正文）}
本题计群行列式的结构分解一个目标。Lemma 4.8 仅为支撑，不另计数。标准先修是 Maschke 定理、复群代数的矩阵块分解和正规表示中不可约表示的重数。题目仍须把群指标矩阵识别为正规表示算子，再通过矩阵块的独立坐标判定因子不可约性及指数；没有把群行列式分解本身作为先修。

所用两页均实际查看原页图像；原命题、证明和变量约定保留，\(Id\)、\(\det|_{V_i}\) 按原文记法。只规范 TeX 空格、换行与公式环境。没有保存下载 PDF 原件，\texttt{sources/mit-representation/014\_assistant\_transcription.tex} 是助手转录。
\subsection*{可修改的简短标注（非作者正文）}
$c$：群指标多项式的不可约分解

$a$：正规表示；群代数矩阵块分解；Maschke 定理；行列式；独立线性坐标；齐次多项式

\texttt{cross\_theory\_status = yes}。把群指标矩阵转换为正规表示的算子值多项式，借不可约表示的矩阵块分解确定多项式因子及其幂次。位置：Theorem 4.7 全证，PDF第3页。

全部标注均为非穷尽、可修改初稿；未标注不表示未使用。
\subsection*{许可与修改记录}
原文来自 MIT OpenCourseWare，作者与课程如上。按来源网站标示的 Creative Commons Attribution--NonCommercial--ShareAlike 4.0 许可复用；本转录及整理沿用同一许可。2026-09-28：助手转录、加入来源定位与简短标注；不暗示作者或 MIT 认可本次整理。许可：\url{https://creativecommons.org/licenses/by-nc-sa/4.0/}。
\end{document}
